\section*{Chapter \ref{ch:s1:short-interest-borrow-and-crowding} --- Short Interest, Borrow and Crowding}

\begin{solution}{exo:s1:short-interest-borrow-and-crowding:1}
$0.1 \times 80 = 8$ million shares short against 0.8 million traded a day: 10 days to cover.
\end{solution}

\begin{solution}{exo:s1:short-interest-borrow-and-crowding:2}
A lends a share to B, who sells it short to C; C lends the same share to D, who sells it short to E. One share has been sold short twice, and short interest
counts two. Repeat and short interest passes the float, as the SEC's staff explained for GameStop.
\end{solution}

\begin{solution}{exo:s1:short-interest-borrow-and-crowding:3}
At 80\% utilisation the fee is $0.25\% + 39.75\% \times 0.5^3 = 5.22\%$ a year. On 5\% of capital for nine months: $0.05 \times 5.22\% \times 0.75 = 0.20\%$ of
capital.
\end{solution}

\begin{solution}{exo:s1:short-interest-borrow-and-crowding:4}
The fee is highest where utilisation is highest, which is a narrower and more expensive set of names than the most shorted as a share of shares: the fee
book's short side pays 7.9\% a year on average against the short-interest book's 5.4\%. The fee book's names underperform by less than their fees, and its
signal is weaker (IC 0.006 against 0.061), because the fee is flat for most stocks and then saturates.
\end{solution}

\begin{solution}{exo:s1:short-interest-borrow-and-crowding:5}
It pays no fees, faces no recalls or squeezes, and needs only that the most-shorted stocks do worse than the rest on average; the long--short version must
also short the names where shorting is most expensive and most dangerous.
\end{solution}

\begin{solution}{exo:s1:short-interest-borrow-and-crowding:6}
Each position is $0.5/20 = 2.5\%$ of capital; a twentyfold rise is a gain of 19 times for the stock, so the book loses $2.5\% \times 19 = 47.5\%$ of capital.
\end{solution}

\begin{solution}{exo:s1:short-interest-borrow-and-crowding:7}
The median losses are 2.8\% (a quarter covered in five days), 4.0\% (half in five days) and 5.7\% (half in ten days). Under the square-root law the daily
move grows with the square root of the covered share; spreading the same buying over more days lowers each day's move but, in the chapter's stress, the
price keeps rising for longer, so the cumulative move grows. Real squeezes are worse: other buyers join, and short sellers facing losses cover faster.
\end{solution}

\begin{solution}{exo:s1:short-interest-borrow-and-crowding:8}
Short interest is known only when FINRA publishes it, days after the settlement date; using it at the settlement date is look-ahead. The backtest must also
charge borrow fees and recalls on its shorts, and report equal- and value-weighted results, since the effect is weak in value-weighted portfolios.
\end{solution}

\begin{solution}{pb:s1:short-interest-borrow-and-crowding:1}
\begin{enumerate}
\item 122.97\% of float in January 2021; when shares sold short are lent again by their buyers.
\item From a close of \$19.95 on January 12 to \$347.51 on January 27 (more than 1\,600\% above January 11), an intraday high of \$483.00 on January 28, and
about 2\,700\% from the January 8 low to that high.
\item Covering by major short sellers coincided with some sharp rises, but it was a small fraction of buying and prices stayed high after it waned.
\item Twice a month under FINRA Rule 4560, due two business days after the settlement date and published later: weeks old when used.
\item Heavily shorted stocks underperformed by a risk-adjusted 1.16\% over 20 days (2000--2004); many patterns are not robust, and value-weighted high
short-interest portfolios did not underperform.
\item Lendable supply around 20\% of shares; informed shorts following the planted expected return and uninformed noise; fees rising above 60\%
utilisation to 40\%; recalls more likely at high utilisation.
\item 0.061, 0.048, 0.006 and 0.058.
\item 8.3\% a year gross, 2.7\% of fees and 0.95\% of costs; Sharpe ratios 3.33, 2.24 and 1.84.
\item Its names underperform by less than their fees (7.9\% a year on the short side).
\item 0.94\% a year over the equal-weighted market, $t = 7.2$, with no fees.
\item Stocks with more short-selling risk (fees that can jump, loans that can be recalled) have lower returns, less price efficiency and less short
selling.
\item A daily probability rising from 0.05\% to 2\% with utilisation above 80\%: high utilisation is where heavily shorted names are.
\item \textbf{Named result.} The short-interest decile book earns 8.3\% a year before borrow fees and 4.6\% after fees and costs (Sharpe ratio 3.33, 2.24
after fees); a book short the twenty most crowded names with half its capital loses a median 4.0\% of capital, and at worst 8.3\%, if half their short
interest covers within five days, against a worst actual five-day loss of 3.0\%.
\item The square-root law is fitted to ordinary orders, not to buying days of volume with other buyers joining.
\item It costs 47.5\% of capital.
\item Cap each name by days to cover and the book's total crowding exposure, and stress it with squeezes rather than history.
\item Avoidance.
\item Daily securities-lending data (utilisation, fees, recalls) and daily short-sale volume, beyond the twice-monthly short interest.
\item The tail of a short book is outside most histories; it has to be imposed by stress scenarios.
\item The borrow fee, the recall risk and an unbounded loss.
\end{enumerate}
\end{solution}

\begin{solution}{iq:s1:short-interest-borrow-and-crowding:1}
That informed traders expect it to fall, on average, and that it is expensive and risky to short; high short interest also means potential forced buying
if short sellers must cover.
\end{solution}

\begin{solution}{iq:s1:short-interest-borrow-and-crowding:2}
Borrow fees that can rise, loans that can be recalled, margin calls, unbounded losses, and squeezes when many short sellers must buy at once.
\end{solution}

\begin{solution}{iq:s1:short-interest-borrow-and-crowding:3}
Score each short on short interest, utilisation and days to cover; limit positions by days to cover and aggregate the book's crowded exposure; stress with
squeeze scenarios; watch fees and recalls daily.
\end{solution}

\begin{solution}{iq:s1:short-interest-borrow-and-crowding:4}
Because shorting it costs a fee that can exceed the underperformance, and carries recall and \omterm{def:s1:short-interest-borrow-and-crowding:squeeze}{squeeze risk} that eats the average; the signal is better used
to avoid the names.
\end{solution}

\begin{solution}{iq:s1:short-interest-borrow-and-crowding:5}
The settlement date of each short-interest report and its publication date; fee and utilisation data dated when observed; volume windows aligned with them.
Without them the backtest uses information before it was public.
\end{solution}

\begin{solution}{iq:s1:short-interest-borrow-and-crowding:6}
Each day's move is about $\eta\sigma\sqrt{Q/(dV)}$, so over $d$ days the price rises by about $d\,\eta\sigma\sqrt{Q/(dV)} = \eta\sigma\sqrt{dQ/V}$ if the
impact persists. It understates a squeeze because the law is estimated on orders far smaller than daily volume, and because a squeeze feeds on itself:
rising prices force more covering and attract momentum buyers, which the law does not contain.
\end{solution}
